The D1 mini board\+: \href{http://wemos.cc}{\tt http\+://wemos.\+cc}

Tool used to convert .bmp to byte arrays\+: \href{http://javl.github.io/image2cpp/}{\tt http\+://javl.\+github.\+io/image2cpp/}

\subsection*{Materials\+:}


\begin{DoxyItemize}
\item We\+Mos D1 Mini
\item Nokia 5110 P\+C\+D8544 84x48 monochrome L\+CD
\item 8 jumper wires
\item breadboard
\end{DoxyItemize}

\subsection*{Dependencies\+:}

Add these to your Arduino I\+DE using Manage Libraries.


\begin{DoxyItemize}
\item \href{https://github.com/adafruit/Adafruit-GFX-Library}{\tt Adafruit G\+FX Library}
\item \href{https://github.com/adafruit/Adafruit-PCD8544-Nokia-5110-LCD-library}{\tt Adafruit P\+C\+D8544 Nokia 5110 L\+CD library}
\begin{DoxyItemize}
\item There is a pull request that adds E\+S\+P8266 support\+:
\item \href{https://github.com/adafruit/Adafruit-PCD8544-Nokia-5110-LCD-library/pull/27}{\tt adafruit/\+Adafruit-\/\+P\+C\+D8544-\/\+Nokia-\/5110-\/\+L\+C\+D-\/library\#27}
\item If this has not been merged yet, you will need to manually copy the changes across
\end{DoxyItemize}
\end{DoxyItemize}

\subsection*{Breadboard connections\+:}

\tabulinesep=1mm
\begin{longtabu} spread 0pt [c]{*3{|X[-1]}|}
\hline
\rowcolor{\tableheadbgcolor}{\bf We\+Mos D1 Mini }&{\bf Nokia 5110 P\+C\+D8544 L\+CD }&{\bf Description  }\\\cline{1-3}
\endfirsthead
\hline
\endfoot
\hline
\rowcolor{\tableheadbgcolor}{\bf We\+Mos D1 Mini }&{\bf Nokia 5110 P\+C\+D8544 L\+CD }&{\bf Description  }\\\cline{1-3}
\endhead
D2 (G\+P\+I\+O4) &0 R\+ST &Output from E\+SP to reset display \\\cline{1-3}
D1 (G\+P\+I\+O5) &1 CE &Output from E\+SP to chip select/enable display \\\cline{1-3}
D6 (G\+P\+I\+O12) &2 DC &Output from display data/command to E\+SP \\\cline{1-3}
D7 (G\+P\+I\+O13) &3 Din &Output from E\+SP S\+PI M\+O\+SI to display data input \\\cline{1-3}
D5 (G\+P\+I\+O14) &4 Clk &Output from E\+SP S\+PI clock \\\cline{1-3}
3\+V3 &5 Vcc &3.\+3V from E\+SP to display \\\cline{1-3}
D0 (G\+P\+I\+O16) &6 BL &3.\+3V to turn backlight on, or P\+WM \\\cline{1-3}
G &7 Gnd &Ground \\\cline{1-3}
\end{longtabu}
These connections are common to all examples.

The backlight (BL) is optional. If you would like full brightness all the time, connect BL to 3\+V3 and save D0 for something else.

\subsection*{Byte Arrays}

I used this tool to convert bitmap images into byte arrays. \href{http://javl.github.io/image2cpp/}{\tt http\+://javl.\+github.\+io/image2cpp/}

The Adafruit library {\ttfamily draw\+Bitmap()} expects steam of bytes which are drawn horizontally. The first {\ttfamily 0x\+FF} in the array represents the first 8x pixels in the top left corner moving to the right.

Each rows last byte needs to be padded with least significant bits to make it a full byte.

I expected a 5x2 black square image would be exactly 10 bits in the array, however, it appears you need 16 bits with 3 padding bits at the end of each row.

The padding least significant bits in the last byte per row are seemingly ignored.


\begin{DoxyCode}
1 // 5x2 black rectangle
2 ABCDE
3 FGHIJ
4 
5 // last byte on each row zero padded
6 ABCDE000
7 FGHIJ000
8 
9 // as bits
10 11111000
11 11111000
12 
13 // as bytes
14 0B11111000, 0B11111000
15 0xF8, 0xF8
\end{DoxyCode}



\begin{DoxyCode}
1 // 9x2 black rectangle
2 ABCDEFGHI
3 JKLMNOPQR
4 
5 // last byte on each row zero padded
6 ABCDEFGH I0000000
7 JKLMNOPQ R0000000
8 
9 // as bits
10 11111111 10000000
11 11111111 10000000
12 
13 // as bytes
14 0B11111111, 0B10000000, 0B11111111, 0B10000000
15 0XFF, 0x80, 0XFF, 0x80
\end{DoxyCode}


To make the image2cpp tool output an Adafruit compatible byte array, you need to upload an image then round up the width to make it divisble by 8. eg. Upload a 84x48 image, then adjust the width to 88x48 then generate the bytes. Upload a 9x2 image and you need to change the size to 16x2 then generate.

The internal buffer in the library stores the bits in 6 rows of 84x columns of 8 bits. Completely different to what {\ttfamily drap\+Bitmap()} accepts! 